\begin{table}[htbp]
\centering
\caption{典型案例下单阶段检索与双阶段检索结果对比}
\label{tab:retrieval_case_compare}
\footnotesize
\setlength{\tabcolsep}{4pt}
\renewcommand{\arraystretch}{1.25}
\begin{tabular}{p{0.17\linewidth} p{0.31\linewidth} p{0.31\linewidth} p{0.14\linewidth}}
\hline
查询内容 & 单阶段检索结果摘要（Top-1/Top-5） & 双阶段检索结果摘要（Top-1/Top-5） & 对比结论 \\
\hline
试卷命题与前两届试题的重复率限制是多少？ &
Top-1 为附件1-9第3个切块（第1页），内容是“ A、B 两套试卷中重复题目不得超过30\% ”；正确条款附件1-9第2个切块（第1页）仅排第2。 &
Top-1 调整为附件1-9第2个切块（第1页），直接命中“试题与前两届试题不得重复超过30\%”；附件1-9第3个切块退居次位。 &
双阶段检索纠正了“相邻条款替代真实证据”的问题，使首位结果具备直接作答性。 \\
\hline
开卷考试中直接从教材找答案的题目分值限制是多少？ &
Top-1 为附件1-7第11个切块（第6页），对应教材编写绩效“最高不超过600分”；Top-2 与 Top-3 仍为教学任务计分与绩效条款；正确证据附件1-9第2个切块（第1页）仅排第4。 &
Top-1 调整为附件1-9第2个切块（第1页），直接给出“可从教材中找到答案的题目分值不应超过总分值的20\%”；Top-2 为同一细则中的相邻命题条款。 &
双阶段检索显著压制了由“教材”“分值”等表层特征引入的噪声，首条结果由无关计分规则转为命题约束条款。 \\
\hline
毕业要求和课程体系的合理性评价周期是几年？ &
Top-1 为附件1-14第3个切块（第1页），给出毕业要求达成度评价“四年一次”的近似条款；Top-2 为附件1-5第6个切块（第2页），给出课程体系“四年修订、每年局调”；真正关键的附件1-15第7个切块（第2页）仅排第3，附件1-15第14个切块（第3页）在单阶段排第6。 &
Top-1 调整为附件1-15第7个切块（第2页），明确“毕业要求至少每二年综合评价一次”；同时附件1-15第14个切块（第3页）升至第4，使“课程体系每两年局部调整”的证据进入前列。 &
双阶段检索改善了跨块、跨页证据的聚合顺序，使“二年”这一最终结论具备更完整的支撑链。 \\
\hline
\end{tabular}
\end{table}
